\documentclass[10pt,a4paper]{amsart}
\usepackage[margin=19mm]{geometry}
\usepackage{amsmath,amssymb,mathtools}
\usepackage[scr=rsfs,cal=euler]{mathalfa}
\usepackage{xfrac}
\usepackage[unicode,hidelinks,bookmarks=false]{hyperref}
\newcommand{\ie}{i.e.\ }
\newcommand{\cf}{cf.\ }
\newcommand{\resp}{respectively\ }
\newcommand{\Subsection}{\subsection}
\providecommand{\nobreakdash}{\nobreak\hbox{-}}
\setlength{\emergencystretch}{2em}
\multlinegap0pt
\usepackage{bm}
\usepackage{bbm}
\usepackage{dsfont}
\usepackage{xspace}
\usepackage{cancel}
\usepackage{mathtools}
\usepackage{amsthm}
\makeatletter
\@ifundefined{definition}{\newtheorem{definition}{Definition}}{}
\@ifundefined{corollary}{\newtheorem{corollary}[definition]{Corollary}}{}
\@ifundefined{example}{\newtheorem{example}[definition]{Example}}{}
\@ifundefined{proposition}{\newtheorem{proposition}[definition]{Proposition}}{}
\@ifundefined{theorem}{\newtheorem{theorem}[definition]{Theorem}}{}
\makeatother
\newcommand{\Adjoint}{\operatorname{Ad}}
\newcommand{\R}{\mathbb{R}}
\newcommand{\coadjoint}{\operatorname{Ad}^*}
\newcommand{\fancyg}{\mathfrak{g}}
\newcommand{\rank}{\operatorname{rank}}
\newcommand{\rdbracket}{[\mathbb{R}^d]}
\title[Fourier Inversion on the Group of Signatures]{Fourier Inversion on the Group of Signatures}
\author{Frank Filbir, Davide Nobile, Marco Rauscher}
\date{Source: arXiv:2509.06106v1 (CC BY 4.0)}
\begin{document}
\maketitle

\noindent\emph{Source and licence.} Fourier Inversion on the Group of Signatures. Version arXiv:2509.06106v1 (CC BY 4.0), distributed under the Creative Commons Attribution 4.0 International licence, as recorded in the arXiv licence metadata for this exact version and shown on the abstract page https://arxiv.org/abs/2509.06106v1. Full rights notice: licenses/arxiv-2509.06106-CC-BY-4.0.txt.\par\smallskip

\noindent\textit{Editorial scope.} Counted unit: the Theorem whose source label is \texttt{thm: characterization of gen orbits} in the article (source0.tex lines 731--743, source label \texttt{thm: characterization of gen orbits}), delivered together with the article's own complete original proof of it (source0.tex lines 744--814, one \texttt{proof} environment of 71 lines). No mathematics is written, repaired or re-derived: statement and proof are the article's own text line for line. Required context retained uncounted: ex: explicit orbits of g2 (lines 437--465); eq: condition for orbit in g2 (lines 454--456); eq: conditions N2 (lines 458--460); ex: explicit orbit of g3 (lines 466--490); eq: necessary condition orbit in g\_3 (lines 484--486); ex: general orbits of g2 pt1 (lines 521--540); prop: existence of invertible matrix (lines 541--549); cor: max rank of blm (lines 656--676); ex: general orbits of g2 pt2 (lines 693--695); ex: general orbits of g\_3 (lines 696--726); eq: ex. space spanned g\_3 (lines 717--719). Every definition, display and label of those lines is retained unchanged. Omitted from this package and not redistributed: 1--436, 457--457, 461--465, 487--520, 550--655, 677--692, 720--730, 815--1284. Nothing inside a retained excerpt is omitted or altered: every retained body is delivered exactly as the article's lines are written, so no omission receipt of any kind accompanies this package. Citations: the retained excerpts carry no citation command at all, so no bibliography entry of the article is transcribed and no reference is redistributed. Reference adaptation: none was needed and none was made; every reference inside the delivered unit resolves inside it, so no unresolved reference or citation is produced. Printed slips kept literal: no printed word, symbol, hypothesis or number of the counted statement or of its proof was altered. Layout and class: the article's own class is \texttt{sn-jnl} with packages including graphicx, multirow, amsmath, amssymb, amsthm, amsfonts, mathrsfs, appendix, xcolor, textcomp, manyfoot, booktabs, algorithm, algorithmicx; the wrapper is the lane's pinned \texttt{amsart} editorial wrapper at 10pt on A4, which restates the theorem-like environments the retained text needs and adds the article's own macro definitions that the retained text needs (\texttt{\textbackslash Adjoint}, \texttt{\textbackslash R}, \texttt{\textbackslash coadjoint}, \texttt{\textbackslash fancyg}, \texttt{\textbackslash rank}, \texttt{\textbackslash rdbracket}), with extra wrapper packages (bm, bbm, dsfont, xspace, cancel, mathtools). The delivered numbering is the unit's own. Assets: no retained span contains a figure environment, a graphics-inclusion command, a PDF-page inclusion, a table or a TikZ picture; no image file is copied into this package, the package declares no asset dependency and no image file is redistributed with it. Licence: version arXiv:2509.06106v1 of this article is distributed under the Creative Commons Attribution 4.0 International licence, as recorded in the arXiv licence metadata for this exact version and shown on the abstract page https://arxiv.org/abs/2509.06106v1. A full rights notice is delivered at licenses/arxiv-2509.06106-CC-BY-4.0.txt. Characters: every retained excerpt is pure ASCII, so the delivered engine, XeLaTeX with its default Latin Modern text font, reports no missing character.
	\begin{example}\label{ex: explicit orbits of g2}
		Let $\{X_1, \ldots, X_d\}$ be an orthonormal basis of $\mathbb{R}^d$ and  set $X_{ij} = [X_i,X_j]$. Then, $\{X_{ij}\}_{1\leq i<j\leq d}$ is a basis of $[\mathbb{R}^d, \mathbb{R}^d]$ so that $\{X_1, \ldots, X_d\} \cup \{X_{ij}\}_{1\leq i<j\leq d}$ is a basis of $\fancyg_2$. Let  $\{\ell_1, \ldots, \ell_{d}\}$ and $\{\ell_{ij}\}_{1\leq i<j\leq d}$ be the corresponding dual bases.\\
		Let $\ell = \sum_i \alpha_i \ell_i + \sum_{i,j} \beta_{i,j} \ell_{i,j} \in \fancyg_2^*$ and take an arbitrary element $X = \sum_{i} g_{i}X_i +  \sum_{i,j} g_{i,j} X_{i,j} \in \fancyg_2$, where $\alpha_i, \beta_{i,j}, g_{i},g_{i,j} \in \mathbb{R}$ for all $i,j$. Then, by definition of the adjoint action, we have that
		\begin{align*}
			&\Adjoint(\operatorname{exp}-X)X_i =X_i+\sum_{j=1}^d g_j[X_i, X_j] = X_i + \sum_{j=1}^d g_j X_{i,j}\\ %\quad \text{and}\\
			& \Adjoint(\operatorname{exp}-X)X_{i,j} = X_{i,j} \;.
		\end{align*}
		To compute the orbit of $\ell$ we need to check how the functional  $(\coadjoint \exp X)\ell$ operates on the basis elements of $\fancyg_2$. By the definition of coadjoint action and the previous equations we have that
		\begin{align*}
			&((\coadjoint \exp X)\ell)(X_i) = \ell(\Adjoint(\exp -X)X_i) = \ell\left(X_i + \sum_{j=1}^d g_j X_{i,j}\right) = \alpha_i + \sum_{j=1}^d g_j\beta_{i,j} \\
			&((\coadjoint \exp X)\ell)(X_{i,j}) = \ell(\Adjoint(\exp -X)X_{i,j}) = \ell(X_{i,j}) = \beta_{i,j}
		\end{align*}
		Therefore, writing $(\coadjoint \exp X)\ell$ using the dual basis of $\fancyg_2^*$ we obtain that the orbit of $\ell$ is given by
		\begin{align*}
			\mathcal{O}_{\ell} = (\coadjoint G) \ell &= \left\{\sum_i(\alpha_i+\sum_jg_{j}\beta_{i,j})\ell_i+\sum_{i,j}\beta_{i,j}\ell_{i,j} \;:\; (g_{1}, \ldots, g_d) \in \mathbb{R}^d\right\}
		\end{align*}
		From this formula we see that a functional $\tilde{\ell} = \sum_i \tilde{\alpha}_i {\ell}_i + \sum_{i,j} \tilde{\beta}_{i,j} {\ell}_{i,j} \in {\fancyg}_2^*$ is in the same orbit as $\ell$ if and only if, $\tilde{\beta}_{i,j}= \beta_{i,j}$ for all $i,j = 1, \ldots, d$ {and} there are $g_1, \ldots, g_d \in \R$ such that
		\begin{align}\label{eq: condition for orbit in g2}
			\tilde{\alpha}_i = \alpha_i + \sum_{j = 1}^dg_j \beta_{i,j} \quad \forall i= 1, \ldots, d 
		\end{align}
		Note that if we define the matrix $B \in \mathbb{R}^{d\times d}$ by $B_{i,j} = \beta_{i,j}$, and $\alpha = (\alpha_1, \ldots, \alpha_d)$, $\tilde{\alpha} = (\tilde{\alpha}_1, \ldots, \tilde{\alpha}_d)$, then (\ref{eq: condition for orbit in g2}) can be rewritten as
		\begin{align}
			B\cdot g = \tilde{\alpha}-\alpha \quad \text{for some } g\in \mathbb{R}^d \;. \label{eq: conditions N2}
		\end{align}
		In particular, if the matrix $B$ is invertible, condition ($\ref{eq: conditions N2}$) is satisfied for any $\alpha, \tilde{\alpha} \in \mathbb{R}^d$. Hence, in this case, the orbit of $\ell$ is a $d$-dimensional submanifold of $\fancyg_2^*$ and it is simply given by 
		\begin{equation*}
			(\coadjoint G) \ell = \sum_{i,j} \beta_{i,j} \ell_{i,j} + \left\{\sum_{i=1}^d \alpha_i \ell_i \; | \; \alpha_1, \ldots, \alpha_d \in \mathbb{R}\right\} = \sum_{i,j} \beta_{i,j} \ell_{i,j} + \R\text{-span }\{\ell_1, \ldots, \ell_d\} \;.
		\end{equation*}
	\end{example}

	\begin{example}\label{ex: explicit orbit of g3}
		We now consider the group $G_3 = G_3(\R^d)$.\\
		Let $\{X_1^1, \ldots, X_d^1\}$, $\{X_1^2, \ldots, X_{m_2}^2\}$, $\{X_1^3, \ldots, X_{m_3}^3\}$ be bases of $\R^d$, $\rdbracket^2$ and $\rdbracket^3$ respectively. 
		Let $\ell = \sum_{i,k} \alpha_i^k\ell_i^k \in \fancyg_3^*$ and take an arbitrary $X = \sum_{i,k} g_i^kX_i^k$, then
		\begin{align*}
			\Adjoint(\exp -X)X_i^1 &= X_i^1 + [X_i^1, X] +\frac{1}{2}[X,[X_i^1,X]]\\
			&= X_i^1 + \sum_{j=1}^{d} g_j^1[X_i^1,X_j^1] + \sum_{j=1}^{m_2} g_j^2[X_i^1,X_j^2] + \frac{1}{2} \sum_{j,s = 1}^d g_j^1g_s^1[X_j^1[X_i^1,X_s^1]]\\
			\Adjoint(\exp -X)X_i^2 &= X_i^2 + [X_i^2, X] = X_i^2 + \sum_{j=1}^{d} g_j^1[X_i^2,X_j^1] \\
			\Adjoint(\exp -X)X_i^3 &= X_i^3 \;.
		\end{align*}
		Hence, using the dual basis as in the previous example, we obtain that the orbit of $\ell$ is given by
		\begin{align*}
			\mathcal{O}_{\ell}= &\Bigg\{\sum_{i=1}^d \left(\ell(X_i^1) + \ell([X_i^1, X]) +\frac{1}{2}\ell([X,[X_i^1,X]])\right)\ell_i^1\\
			&+\sum_{i=1}^{m_2} \left(\ell(X_i^2) + \ell([X_i^2, X])\right)\ell_i^2 + \sum_{i=1}^{m_3} \alpha_i^3\ell_i^3\; \Big| \; X \in \fancyg_3\Bigg\}\\
			= &\Bigg\{\sum_{i=1}^d \left(\alpha_i^1 + \sum_{j=1}^{d} g_j^1\ell([X_i^1,X_j^1]) + \sum_{j=1}^{m_2} g_j^2\ell([X_i^1,X_j^2]) + \frac{1}{2} \sum_{j,s = 1}^d g_j^1g_s^1\ell([X_j^1[X_i^1,X_s^1]]) \right)\ell_i^1\\
			&+ \sum_{i=1}^{m_2}\left(\alpha_i^2 + \sum_{j=1}^{d} g_j^1\ell([X_i^2,X_j^1])\right)\ell_i^2 + \sum_{i=1}^{m_3} \alpha_i^3\ell_i^3 \;\Big| \; (g_1^1, g_2^1, \ldots, g_{m_2}^2) \in \R^{d+m_2}\Bigg\} \;.
		\end{align*}
		Again, we see that a necessary condition for some $\tilde{\ell} = \sum_{i,k} \tilde{\alpha}_i^k \ell_i^k \in \fancyg_3^*$ to be in the same orbit as $\ell$ is that they coincide on $[\R^3]^3$. Further, there must be some $g = (g_1^1, \ldots, g_d^1) \in \mathbb{R}^d$ such that for all $i = 1, \ldots, m_2$
		\begin{equation}\label{eq: necessary condition orbit in g_3}
			\tilde{\alpha}_i^2 = \alpha_i^2 +  \sum_{j=1}^{d} g_j^1\ell([X_i^2,X_j^1])
		\end{equation}
		As before, setting $\alpha = (\alpha_1, \ldots, \alpha_{m_2})$, $\tilde{\alpha} = (\tilde{\alpha_1}, \ldots, \tilde{\alpha}_{m_2})$ and $B \in \R^{m_2 \times d}$ with $B_{i,j} = \ell([X_i^2, X_j^1])$ we can rewrite (\ref{eq: necessary condition orbit in g_3}) as
		\[B \cdot g = \tilde{\alpha} -\alpha \quad \text{for some } g \in \R^d \;.\]
		Note that this is again only a necessary, but not a sufficient condition for $\tilde{\ell}$ to be in $\mathcal{O}_{\ell}$.
	\end{example}

	\begin{example}\label{ex: general orbits of g2 pt1}
		As in Example \ref{ex: explicit orbits of g2} we consider $\fancyg_2(\R^d)$, with strong Malcev basis $\{X_{1,2}, X_{1,3}, \ldots, X_{d-1,d},X_1, \ldots, X_d\}$ . There, we have seen that the orbit of some $\ell = \sum_i \alpha_i \ell_i + \sum_{i,j} \beta_{i,j} \ell_{i,j} \in \fancyg_2^*$ is given by
		\begin{align*}
			\mathcal{O}_\ell &= \left\{\sum_{i=1}^d(\alpha_i+\sum_{j=1}^dg_{j}\beta_{i,j})\ell_i+\sum_{1\leq i<j\leq d}\beta_{i,j}\ell_{i,j} \;:\; (g_{1}, \ldots, g_d) \in \mathbb{R}^d\right\} \;.
		\end{align*}
		Clearly, the dimension of $\mathcal{O}_\ell$ is smaller or equal $d$. For any $1 \leq m<n\leq d$ we have that $\fancyg_2^*/\fancyg_{2}^*((m, n)) = \R$-span$\{l_{1,2}, \ldots,l_{m,n}\}$ so that
		\begin{align*}
			\dim\left(\mathcal{O}_\ell/\fancyg_{2}^*((m, n))\right)= \dim\Big\{\sum_{\substack{i\leq m,\\j\leq n}}\beta_{i,j}\ell_{i,j}\Big\} = 0 \;. \;
		\end{align*}
		Further, for any $m = 1, \ldots, d$, we have $\fancyg_2^*/\fancyg_{2}^*(m) = \R$-span$\{l_{1,2}, \ldots,l_{d-1,d}, l_1, \ldots, l_m\}$, so that
		\begin{align*}
			\dim\left(\mathcal{O}_\ell/\fancyg_{2}^*(m)\right)&= \dim\left\{\sum_{i=1}^m(\alpha_i+\sum_{j=1}^dg_{j}\beta_{i,j})\ell_i+\sum_{i,j}\beta_{i,j}\ell_{i,j} \;:\; (g_{1}, \ldots, g_d) \in \mathbb{R}^d\right\}\\
			&=\dim\left\{\sum_{i=1}^m\left(\sum_{j=1}^dg_{j}\beta_{i,j}\right)\ell_i \;:\; (g_{1}, \ldots, g_d) \in \mathbb{R}^d\right\}\\
			&=\operatorname{rank} B_\ell(m) \;,
		\end{align*}
		where $B_\ell(m) \in \R^{m\times d}$ is defined by $(B_\ell(m))_{i,j} = \beta_{i,j} = \ell([X_i,X_j])$. By definition, an orbit is in general position if the dimension of $\mathcal{O}_\ell / \mathfrak{g}_2^*(m)$ is maximal for all $m = 1, \ldots, d$. From the previous computation, this condition is equivalent to requiring that the rank of the matrix $B_\ell(m)$ is maximal for each $m$. In particular, if there exists an element $\ell$ such that $B_\ell(d)$ is invertible, then
		\[\dim\left(\mathcal{O}_\ell/\fancyg_{2}^*(m)\right) = m \quad
		\text{for all} \quad m = 1, \ldots, d\]
		and any other $\tilde{\ell} \in \mathfrak{g}_2^*$ is in general position if and only if $B_{\tilde{\ell}}(d)$ is invertible. However, at this point, it is not clear whether such an $\ell$ always exists or, more generally, what the maximal attainable rank of the matrix $B_\ell(d)$ is. The next proposition addresses this problem for arbitrary dimensions.
	\end{example}

	\begin{proposition}\label{prop: existence of invertible matrix}
		Let $\{X_1^N, \ldots, X_{m_N}^N, \ldots,X_1^1, \ldots, X_{m_1}^1\}$ be a strong Malcev basis for $\fancyg_N$ such that for all $k \in \{1, \ldots, N\}$ the elements $X_1^k, \ldots, X_{m_k}^k$ are a basis of the $k$-th layer of $\fancyg_N$. For $k \in \{1,\ldots, \lfloor N/2 \rfloor\}$ define the matrix $B_{\ell}^k \in \mathbb{R}^{m_k\times m_k}$ by
		\[(B_{\ell}^k)_{i,j} = \ell([X_i^k,X_j^{N-k}]) \quad i,j \in \{1, \ldots, m_k\} \;.\]
		Then, there is an $\ell \in \fancyg_N^*$ such that:
		\begin{itemize}
			\item[i)] If $N/2 \notin \mathbb{N}$, or $N/2 \in \mathbb{N}$ with $m_{N/2} \in 2\mathbb{N}$, then $B_{\ell}^k$ is invertible for all $k \leq N/2$.
			\item[ii)] If $N/2 \in  \mathbb{N}$ with $m_{N/2} \in 2\mathbb{N}+1$, then $B_{\ell}^k$ is invertible for all $k<N/2$ and $B_{\ell}^{N/2}$ has rank $m_{N/2}-1$, where the first $m_{N/2}-1$ rows of $B_{\ell}^{N/2}$ are linearly independent.
		\end{itemize}
	\end{proposition}

	\begin{corollary}\label{cor: max rank of blm}
		Let $\{X_1^N, \ldots, X_{m_N}^N, \ldots,X_1^1, \ldots, X_{m_1}^1\}$ be a strong Malcev basis for $\fancyg_N$ such that for all $k \in \{1, \ldots, N\}$ the elements $X_1^k, \ldots, X_{m_k}^k$ are a basis of $k$-th layer of $\fancyg_N$. For $k \in \{1,\ldots, N-1\}$ and $m \in \{1,\ldots,m_{N-k}\}$ define the matrix $B_{\ell}^k(m) \in \mathbb{R}^{m_k\times m}$ by
		\[(B_{\ell}^k(m))_{i,j} = \ell([X_i^k,X_j^{N-k}]) \quad i \in \{1, \ldots, m_k\}, \;j \in \{1, \ldots, m\} \;.\]
		Then, for any $m \in \{1, \ldots, m_{N-k}\}$
		\begin{align*}
			&\operatorname{dim}(k,m) \coloneqq \max_{\ell \in \fancyg_{N}^*}\left(\operatorname{rank}B_{\ell}^k(m)\right)=
			\min\left\{m_k,m_{N-k},m\right\} \text{, if } k \neq N/2\\
			&\operatorname{dim}(N/2,m) \coloneqq \max_{\ell \in \fancyg_{N}^*}\left(\operatorname{rank}B_{\ell}^{N/2}(m)\right) = \begin{cases}
				\min\left\{m_{N/2},m\right\} &\text{ if } m_{N/2} \text{ is even}\\
				\min\left\{m_{N/2}-1,m\right\} &\text{ if } m_{N/2} \text{ is odd}\\
			\end{cases}
		\end{align*}
		and there is an $\ell \in \fancyg_{N}^*$ such that $\operatorname{rank}B_{\ell}^k(m) = \dim(k,m)$ for all $k,m$.\\
		Further, the following hold for any $k < N/2$ :
		\begin{itemize}
			\item[i)] $\operatorname{rank}B_{\ell}^k(m) = \dim(k,m)$ for all $m = 1, \ldots, m_{N-k}$ if and only if $\rank B_{\ell}^k = \dim(k,m_k)$, i.e. if and only if $B_{\ell}^k$ is invertible.
			% \item[ii)] If $k = N/2$ and $m_{N/2} \in 2\mathbb{N}+1$. Then, $\operatorname{rank}B_{\ell}^{N/2}(m) = \dim({N/2},m)$ for all $m = 1, \ldots, m_{N/2}$ if and only if $\rank B_{\ell}^{N/2} = \dim({N/2},m_{N/2}) = m_{N/2}-1$ and the first $m_{N/2}-1$ rows of  $B_{\ell}^{N/2}$ are linearly independent.
			\item[ii)] If $B_{\ell}^k$ is invertible, then $\operatorname{rank}B_{\ell}^{N-k}(m) = \dim(N-k,m)$ for all $m = 1, \ldots, m_{k}$.
		\end{itemize}
		% In particular,  $\operatorname{rank}B_{\ell}^k(m) = \dim(k,m)$ for all $k = 1,\ldots,N$ and $m = 1, \ldots, m_k$, if and only if $\ell$ satisfies $i)$ resp. $ii)$ from Proposition \ref{prop: existence of invertible matrix}.
	\end{corollary}

	\begin{example}\label{ex: general orbits of g2 pt2}
		Combining our observations from Example \ref{ex: general orbits of g2 pt1} and Proposition \ref{prop: existence of invertible matrix} we see that if $d$ is even, then $\ell \in \fancyg_2^*(\R^d)$ is in general position if and only if $B_{\ell}^1$ is invertible. On the other hand, if $d$ is odd, then $\ell$ is in general position if and only if the first $d-1$ rows of $B_{\ell}^1$ are linearly independent.
	\end{example}

	\begin{example}\label{ex: general orbits of g_3}
		Consider the group $G_3 = G_3(\R^d)$ for $d>3$. By example \ref{ex: explicit orbit of g3} we know that the coadjoint orbit of any $\ell = \sum_{i,j}\alpha_i^j\ell_i^j \in \fancyg_3^*$ is given by
		\begin{align*}
			\mathcal{O}_{\ell}= &\Bigg\{\sum_{i=1}^d \left(\alpha_i^1 + \sum_{j=1}^{d} g_j^1\ell([X_i^1,X_j^1]) + \sum_{j=1}^{m_2} g_j^2\ell([X_i^1,X_j^2]) + \frac{1}{2} \sum_{j,s = 1}^d g_j^1g_s^1\ell([X_j^1[X_i^1,X_s^1]]) \right)\ell_i^1\\
			&+ \sum_{i=1}^{m_2}\left(\alpha_i^2 + \sum_{j=1}^{d} g_j^1\ell([X_i^2,X_j^1])\right)\ell_i^2 + \sum_{i=1}^{m_3} \alpha_i^3\ell_i^3 \;\Big| \; (g_1^1, g_2^1, \ldots, g_{m_2}^2) \in \R^{d+m_2}\Bigg\} \;.
		\end{align*}
		Hence, for any $m \in \{1, \ldots,m_2\}$ we have that
		\begin{align*}
			\dim\left(\mathcal{O}_{\ell}/\fancyg_{2}^*(2, m)\right) &=   \dim\Bigg\{\sum_{i=1}^{m}\left(\alpha_i^2 + \sum_{j=1}^{d} g_j^1\ell([X_i^2,X_j^1])\right)\ell_i^2 + \sum_{i=1}^{m_3} \alpha_i^3\ell_i^3 \;\Big| \; (g_1^1,\ldots, g_{d}^1) \in \R^{d}\Bigg\} \\
			&=   \dim\Bigg\{\sum_{i=1}^{m}\left(\sum_{j=1}^{d} g_j^1\ell([X_i^2,X_j^1])\right)\ell_i^2 \;\Big| \; (g_1^1,\ldots, g_{d}^1) \in \R^{d}\Bigg\}\\ 
			&= \operatorname{rank} -B_{\ell}^{1}(m)^T = \operatorname{rank} B_{\ell}^{1}(m) \;.
		\end{align*}
		By Corollary \ref{cor: max rank of blm} the maximal attainable rank of $B_{\ell}^{1}(m)$ is $\min\{d,m_2,m\} = \min\{d,m\}$. Hence, a necessary condition for $\ell$ to be in general position is that $\operatorname{rank} B_{\ell}^{1}(m) = \min\{m,d\}$ for all $m \leq m_2$. By the same corollary this is equivalent to $B_{\ell}^{1} = B_{\ell}^{1}(d)$ being invertible. Let now $m \in \{1, \ldots, d\}$. Then,
		\begin{align*}
			\dim&\left(\mathcal{O}_{\ell}/\fancyg_{2}^*(1, m)\right) \\
            =& \dim \Bigg\{\sum_{i=1}^m \left(\sum_{j=1}^{d} g_j^1\ell([X_i^1,X_j^1]) + \sum_{j=1}^{m_2} g_j^2\ell([X_i^1,X_j^2]) + \frac{1}{2} \sum_{j,s = 1}^d g_j^1g_s^1\ell([X_j^1[X_i^1,X_s^1]]) \right)\ell_i^1\\
			&+ \sum_{i=1}^{m_2}\left(\sum_{j=1}^{d} g_j^1\ell([X_i^2,X_j^1])\right)\ell_i^2 \;\Big| \; (g_1^1, g_2^1, \ldots, g_{m_2}^2) \in \R^{d+m_2}\Bigg\} \;.
		\end{align*}
		By the previous observation, we know that the term
		$\sum_{i=1}^{m_2}\left(\sum_{j=1}^{d} g_j^1\ell([X_i^2,X_j^1])\right)\ell_i^2$
		spans a space of dimension at most $d$, so that clearly $\dim\left(\mathcal{O}_{\ell}/\fancyg_{2}^*(1, m)\right) \leq m +d$. Further, we observe that the dimension of the space spanned by the term
		\begin{equation}\label{eq: ex. space spanned g_3}
			\sum_{i=1}^m \left(\sum_{j=1}^{m_2} g_j^2\ell([X_i^1,X_j^2])\right)\ell_i^1
		\end{equation}
		is precisely the rank of the matrix $-B_{\ell}^2(m)^T$. By Corollary \ref{cor: max rank of blm} we know that if $B_{\ell}^1$ is invertible, then $\rank B_{\ell}^2(m)$ is maximal and equal to $m$ for all $m =1, \ldots, d$. 
		Note that the space spanned by (\ref{eq: ex. space spanned g_3}) only depends on $g_1^2, \ldots, g_m^2$ while the space spanned by the $\ell_i^2$ only depends on $g_1^1,\ldots,g_d^1$. Hence, if $B_{\ell}^1$ is invertible, we have that  $\dim\left(\mathcal{O}_{\ell}/\fancyg_{2}^*(1, m)\right) \geq m +d$. Since $B_{\ell}^1$ being invertible was also a necessary condition for $\ell$ to be in general position, we obtain in total that $\ell$ is in general position if and only if $B_{\ell}^1$ is invertible. In this case, for any $m = 1, \ldots, d$
		\begin{align*}
			\dim\left(\mathcal{O}_{\ell}/\fancyg_{2}^*(2, m)\right) = m \quad \text{and} \quad
			\dim\left(\mathcal{O}_{\ell}/\fancyg_{2}^*(1, m)\right) = d + m \;.
		\end{align*}
	\end{example}

	\begin{theorem}\label{thm: characterization of gen orbits}
		Let $\{X_1^N, \ldots, X_{m_N}^N, \ldots,X_1^1, \ldots, X_{m_1}^1\}$ be a strong Malcev basis for $\fancyg_N \neq \fancyg_3(\R^2)$ such that for all $k \in \{1, \ldots, N\}$ the elements $X_1^k, \ldots, X_{m_k}^k$ are a basis of $k$-th layer of $\fancyg_N$, and let $\ell\in \fancyg_{N}^*$. Then, with the notation of Corollary \ref{cor: max rank of blm}, the following are equivalent.
		\begin{itemize}
			\item[i)] The dimension of $\mathcal{O}_{\ell}$ is maximal in $\fancyg_{N}^*/\fancyg_{N}^*((N-k, m)) = \mathbb{R}$-span$\{\ell_1^N, \ldots, \ell_{m}^{N-k}\}$ for all $k \in \{1, \ldots,N-1\}$ and $m \in \{1, \ldots, m_{N-k}\}$, i.e. $\ell$ is in general position.
			\item[ii)] The rank of the matrix $B_{\ell}^{k}(m)$ is maximal among elements of $\fancyg_{N}^*$ for all $k \in \{1, \ldots,N-1\}$ and $m \in \{1, \ldots, m_{N-k}\}$.
			\item[iii)] The rank of the matrix $B_{\ell}^{k}$ is maximal,  among elements of $\fancyg_{N}^*$ for all $k \in \{1, \ldots,\lfloor N/2\rfloor\}$. I.e. $B_{\ell}^{k} = B_{\ell}^{k}(m_k)$ satisfies property $i)$ resp. $ii)$ from Proposition \ref{prop: existence of invertible matrix}.
		\end{itemize}
		If $\ell$ satisfies any of the conditions above, then for all $k \in \{1, \ldots,N-1\}$
		\begin{align}\label{eq: dimension formual for generic orbits}
			\operatorname{dim}\left(\mathcal{O}_{\ell} / \fancyg_N^*(N - k, m)\right) =
			\sum_{s=1}^{k-1}\operatorname{dim}(s,m_s)+ \operatorname{dim}(k,m)
		\end{align}
	\end{theorem}

	\begin{proof}
		We assume $N > 3$ since we already computed the cases $N=2,3$ explicitly in examples \ref{ex: general orbits of g2 pt2} and \ref{ex: general orbits of g_3} respectively.
		Let $\ell = \sum_{s=1}^{N}\sum_{j=1}^{m_s}\alpha_j^s\ell_j^s \in \fancyg_{N}^*$.\\
		The equivalence $ii) \iff iii)$ is precisely what we showed in Corollary \ref{cor: max rank of blm}.\\   
		We now prove the equivalence $i) \iff ii)$ together with the dimension formula (\ref{eq: dimension formual for generic orbits}) by induction over $k$.\\
		$1)\;$ $k = 1$. Let $X = \sum_{s=1}^{N}\sum_{j=1}^{m_s}g_j^sX_j^s \in \fancyg_{N}$. Then, for all $i \in \{1, \ldots, m_{N-1}\}$
		\begin{align*}
			(\Adjoint \operatorname{exp}-X)X_i^{N-1} &= X_i^{N-1} + \sum_{n=1}^{\infty}\frac{1}{n!}(\operatorname{ad}-X)^n X_i^{N-1}\\
			&=  X_i^{N-1} + [X_i^{N-1}, X] = X_i^{N-1} + \sum_{j=1}^{m_1}g_j^1[X_i^{N-1},X_j^1]
		\end{align*}
		Hence, for any $m \in \{1, \ldots, m_{N-1}\}$ the orbit of $\ell$ in $\mathbb{R}$-span$\{\ell_1^N, \ldots, \ell_{m}^{N-1}\}$ is given by
		\begin{align*}
			\mathcal{O}_{\ell} / \fancyg_N^*(N-1, m) = \Bigg\{ 
			& \sum_{i=1}^{m} \left(\alpha_i^{N-1} + \sum_{j=1}^{m_1} \ell([X_i^{N-1}, X_j^1])g_j^1\right) \ell_i^{N-1} \\
			& + \sum_{i=1}^{m_N} \alpha_i^N \ell_i^N \; \Big| \; (g_1^1, \ldots, g_{m_1}^1) \in \mathbb{R}^{m_1} \Bigg\}
		\end{align*}
		From this equation we see that the dimension of $\mathcal{O}_{\ell} / \fancyg_N^*(N-1, m)$ is precisely the dimension of the image, or in other words the rank, of the $m \times m_1$ matrix with entries $\ell([X_i^{N-1}, X_j^1])$. Note that this matrix is precisely $-(B_{\ell}^{1}(m))^T$. Hence, we obtain that the dimension of $\mathcal{O}_{\ell} / \fancyg_N^*(N-1, m)$ is maximal for all $m$ if and only if the rank of $B_{\ell}^{1}(m)$ is. By Corollary \ref{cor: max rank of blm}, we know that the maximal rank of $B_{\ell}^{1}(m)$ is given by $\dim(1,m)$ and it is attained if and only if $B_{\ell}^1$ is invertible\vspace{5pt}.\\  
		\noindent $2)$ Assume now $k>1$. Let $X = \sum_{s=1}^{N}\sum_{j=1}^{m_s}g_j^sX_j^s \in \fancyg_{N}$, then
		\begin{alignat*}{3}
			(\Adjoint \operatorname{exp} - X)X_i^{N-k} 
			& = X_i^{N-k} 
			& & + \sum_{n=1}^{\infty} \frac{1}{n!} (\operatorname{ad} -X)^n X_i^{N-k} \\
			& = X_i^{N-k} 
			& &+ [X_i^{N-k}, X] 
			+ \sum_{n=2}^{\infty} \frac{1}{n!} (\operatorname{ad} -X)^n X_i^{N-k} \\
			& = X_i^{N-k} 
			&& + \sum_{s=1}^{k} \sum_{j=1}^{m_s} g_j^s [X_i^{N-k}, X_j^s] \\
			&& & + \sum_{n=0}^{\infty} \frac{1}{(n+2)!} (\operatorname{ad} -X)^n [[X_i^{N-k}, X], X] \;.
		\end{alignat*}
		Note that the very last sum only depends on the $g_j^s$ with $s \leq k-1$, since for all $n \in \mathbb{N}_0$ $(\operatorname{ad}-X_j^{s_1})^n [[X_i^{N-k}, X_j^{s_2}], X_j^{s_3}] = 0$ if any of the $s_i \geq k$. Hence, if we set
		\begin{equation*}
			R(g_1^1,g_2^1, \ldots, g_{m_{k-1}}^{k-1}) = \sum_{s=1}^{k-1} \sum_{j=1}^{m_s} g_j^s \; \ell([X_i^{N-k}, X_j^s]) + \sum_{n=0}^{\infty} \frac{1}{(n+2)!}\; \ell\left((\operatorname{ad} -X)^n [[X_i^{N-k}, X], X]\right)
		\end{equation*}
		the orbit of $\ell$ in $\fancyg_N^*(N-k, m)$ has the form
		\begin{align}\label{eq: korbit}
			\mathcal{O}_{\ell} / \fancyg_N^*(N-k, m) = &\Bigg\{ 
			\sum_{i=1}^{m} \left(\alpha_i^{N-k} + \sum_{j=1}^{m_k} g_j^k \ell([X_i^{N-k}, X_j^k]) + R(g_1^1,g_2^1, \ldots, g_{m_{k-1}}^{k-1})\right) \ell_i^{N-k} \notag \\
			& + \sum_{s = 1}^{k-1} \sum_{i=1}^{m_s} r_i^{N-s} \ell_i^{N-s} \; \Big| \; (g_1^1,g_2^1, \ldots, g_{m}^k) \in \mathbb{R}^{m_1+m_2+\ldots + m} \Bigg\}
		\end{align}
		where the $r_i^{N-s}$ are some real coefficients that depend on the $g_i^s$.\\
		We need to differentiate two cases. First, assume $1<k\leq N/2$. Then,
		$\operatorname{dim}(N-s,m) = \min\left\{m_s,m\right\}$ for all $s< k$.\\
		Now, the term
		$\sum_{s = 1}^{k-1} \sum_{i=1}^{m_s} r_i^{N-s} \ell_i^{N-s}$ only depends on $(g_1^1,g_2^1 \ldots, g_{m_{k-1}}^{k-1}) \in \mathbb{R}^{m_1+m_2+\ldots + m_{k-1}}$.
		Since, as we mentioned before, also $R(g_1^1,g_2^1, \ldots, g_{m_{k-1}}^{k-1})$ only depends on $(g_1^1,g_2^1 \ldots, g_{m_{k-1}}^{k-1})$ we can rewrite the orbit as
		
		\[\mathcal{O}_{\ell} / \fancyg_N^*(N - k, m) = W_{\ell}^{k-1} + V_{\ell}^k \quad \text{where}\]
		\begin{alignat*}{3}
			W_{\ell}^{k-1} &=  \Bigg\{
			&& \sum_{i=1}^{m} \left(\alpha_i^{N-k}+ R(g_1^1,g_2^1, \ldots, g_{m_{k-1}}^{k-1})\right) \ell_i^{N-k}\\
			&&& + \sum_{s = 1}^{k-1} \sum_{i=1}^{m_s} r_i^{N-s} \ell_i^{N-s} \; \Big| \; (g_1^1,g_2^1 \ldots, g_{m_{k-1}}^{k-1}) \in \mathbb{R}^{m_1+m_2+\ldots + m_{k-1}} \Bigg\} \quad \text{and} \\
			V_{\ell}^k &= \Bigg\{ 
			&&\sum_{i=1}^{m} \left(\sum_{j=1}^{m_k} g_j^k \ell([X_i^{N-k}, X_j^k])\right) \ell_i^{N-k} \; \Big| \; (g_1^k, \ldots, g_{m_k}^k) \in \mathbb{R}^{m_k} \Bigg\} \;.
		\end{alignat*}
		
		Clearly $\max_{\ell' \in \fancyg_{N}^*}\left(\operatorname{dim}W_{\ell'}^{k-1}\right)\leq m_1+\ldots + m_{k-1}$. On the other hand, by induction assumption, if $\operatorname{rank}B_{\ell}^s(m_s)$ is maximal for all $s<k$, then $$\operatorname{dim}W_{\ell}^{k-1} = \operatorname{dim}(N-1,m_1)+\ldots + \operatorname{dim}(N-k+1,m_{k-1})=m_1+\ldots + m_{k-1} \;.$$
		Further, as before, $\dim(V_{\ell}^k) =\operatorname{rank}(B_{\ell}^k(m))$. Hence, if the rank of $B_{\ell}^k(m)$ is also maximal we obtain that
		\begin{align*}
			\operatorname{dim}\left(\mathcal{O}_{\ell} / \fancyg_N^*(N - k, m)\right)
			&= \sum_{s=1}^{k-1}\operatorname{dim}(N-s,m_s)+ \operatorname{dim}(N-k,m)\\
			&= \sum_{s=1}^{k-1}m_s+ \max_{\ell' \in \fancyg_{N}^*}\left(\rank B_{\ell'}^k(m)\right)\\
			&=\max_{\ell' \in \fancyg_{N}^*}\left(\dim W_{\ell'}^{k-1} + \dim V_{\ell}^k\right) = \max_{\ell' \in \fancyg_{N}^*}\big(\operatorname{dim}\left(\mathcal{O}_{\ell'} / \fancyg_N^*(N - k, m)\right) \big)
		\end{align*}
		Conversely, if the dimension of $\mathcal{O}_{\ell} / \fancyg_N^*(N - k, m)$ is maximal, then so is $\dim(V_{\ell}^k) = \operatorname{rank}(B_{\ell}^k(m))$\vspace{5pt}.\\
		We now assume $k>N/2$. \\
		Then, $N-k<N/2$ and $m_k>m_{N-k}$. If $\ell$ is in general position, then we know by induction that $B_{\ell}^{N-k}$ must have maximal rank. By Corollary \ref{cor: max rank of blm} $ii)$ this implies that $B_{\ell}^{k}(m)$ has maximal rank for all $m$, i.e. $\rank B_{\ell}^{k}(m) = \dim(k,m)$.\\
		Conversely, assume that $\operatorname{rank}B_{\ell}^{N-s}(m)=\dim(N-s,m)$ for all $s$. By induction the term $ \sum_{s = 1}^{k-1} \left(\sum_{i=1}^{m_s} r_i^{N-s} \ell_i^{N-s} \right)$ in equation (\ref{eq: korbit}) spans a subspace of dimension $\sum_{s =1}^{k-1}\dim(N-s,m_s)$. Further, for all $m \in \{1, \ldots, m_{N-k}\}$ the term
		\[\sum_{i=1}^{m} \left(\alpha_i^{N-k} + \sum_{j=1}^{m_k} g_j^k \ell([X_i^{N-k}, X_j^k]) + R(g_1^1,g_2^1, \ldots, g_{m_{k-1}}^{k-1})\right) \ell_i^{N-k}\]
		spans a subspace of dimension less or equal $m =\min\left\{m,m_{N-k}\right\} =\min\left\{m,m_{N-k},m_k\right\} = \dim(N-k,m)$ and again equality holds if $B_{\ell}^{k}(m)$ has rank $m$.
		This concludes the proof of $i) \iff ii)$ and of the dimension formula.
	\end{proof}

\end{document}
